\section*{Devoir maison : Ranger les Lagonnet (Tri par Insertion)}

Les quatre frères Lagonnet sont en file indienne, mais ils sont dans le désordre. Le gardien de prison veut les ranger \textbf{du plus petit au plus grand} (ordre croissant).

\subsection*{1. Le principe : "Chacun son tour à sa place"}
On laisse le premier Lagonnet là où il est. Ensuite, on prend le deuxième, puis le troisième, etc. Pour chaque Lagonnet qu'on "sort" de la file :
\begin{itemize}
    \item On regarde son voisin de gauche.
    \item Tant que le voisin de gauche est \textbf{plus grand} que lui, on fait reculer ce voisin d'un cran vers la droite pour laisser une place.
    \item Dès qu'on trouve un voisin plus petit (ou qu'on arrive au début de la file), on insère notre Lagonnet dans l'espace vide.
\end{itemize}

\begin{lstlisting}[language=Python]
def tri_insertion(tab):
    for i in range(1, len(tab)):
        cle = tab[i]          # Le Lagonnet qu'on veut ranger
        j = i - 1             # Son voisin de gauche
        
        # Tant que le voisin est plus grand que le Lagonnet "cle"
        while j >= 0 and tab[j] > cle:
            tab[j + 1] = tab[j] # Le voisin recule d'une place
            j = j - 1           # On regarde le voisin suivant a gauche
            
        tab[j + 1] = cle      # On insere le Lagonnet a sa place
\end{lstlisting}

\begin{Exercise}[title={Analyse et Simulation}]

\Question \textbf{Simulation guidée.} 
Les tailles des Lagonnet sont : \texttt{tab = [1.5, 0.9, 2.1, 1.2]}.
\begin{itemize}
    \item \textbf{Tour 1 ($i=1$)} : La \texttt{cle} est \textbf{0.9} (Petit-Îlot). Son voisin de gauche est \textbf{1.5}.
    \begin{itemize}
        \item Est-ce que $1.5 > 0.9$ ? ............
        \item On décale le 1.5 à droite. Le tableau devient provisoirement \texttt{[1.5, 1.5, 2.1, 1.2]}.
        \item On pose le 0.9 au début. Résultat : \texttt{[0.9, 1.5, 2.1, 1.2]}.
    \end{itemize}
    \item \textbf{Tour 2 ($i=2$)} : La \texttt{cle} est \textbf{2.1} (Mont-Choungui). Son voisin est \textbf{1.5}.
    \begin{itemize}
        \item Est-ce que $1.5 > 2.1$ ? ............
        \item Rien ne bouge, Mont-Choungui est déjà plus grand que son voisin. Résultat : \texttt{[0.9, 1.5, 2.1, 1.2]}.
    \end{itemize}
\end{itemize}

\Question \textbf{À vous de jouer.} 
Terminez le tri pour le dernier tour (\textbf{Tour 3, $i=3$}). La \texttt{cle} est \textbf{1.2}.
\begin{itemize}
    \item Avec qui compare-t-on 1.2 en premier ? .........................
    \item Est-ce qu'on doit décaler des valeurs ? Si oui, lesquelles ? ............................................
    \item Quel est le tableau final trié ? \texttt{[ \quad , \quad , \quad , \quad ]}
\end{itemize}

\Question \textbf{Réflexion sur la difficulté (Complexité).}
\begin{enumerate}
    \item \textbf{Meilleur cas :} Si les Lagonnet arrivent déjà rangés \texttt{[0.9, 1.2, 1.5, 2.1]}, est-ce que la boucle \texttt{while} (le décalage) va s'exécuter souvent ? .........................
    \item \textbf{Pire cas :} Si les Lagonnet arrivent dans l'ordre inverse \texttt{[2.1, 1.5, 1.2, 0.9]}, que doit faire l'ordinateur pour chaque nouveau Lagonnet ? \\ ..................................................................................................................
    \item \textbf{Conclusion :} Si on a $n$ Lagonnet, dans le pire des cas, le temps de calcul augmente comme le carré de $n$ ($n \times n$). Pourquoi dit-on que cet algorithme est "lent" pour les très grandes listes ? \\ ..................................................................................................................
\end{enumerate}

\end{Exercise}